\def\shorttitle{Samper：有序单纯复形的拟拟阵类}
\documentclass[UTF8,zihao=-4,fontset=windows,a4paper]{ctexart}
\usepackage[margin=25mm,headheight=16pt]{geometry}
\usepackage{amsmath,amssymb,amsthm,mathtools,enumitem,fancyhdr,booktabs,longtable,array,xcolor}
\usepackage[unicode,colorlinks,linkcolor=blue!45!black,urlcolor=blue!45!black,bookmarksnumbered]{hyperref}
\setlength{\parskip}{.35em}
\setlength{\emergencystretch}{2em}
\clubpenalty=10000
\widowpenalty=10000
\displaywidowpenalty=10000
\pdfstringdefDisableCommands{\def\quad{ }}
\linespread{1.15}
\setlist{nosep,leftmargin=2em,topsep=.4em}
\pagestyle{fancy}\fancyhf{}\fancyhead[L]{\small 拟阵与单纯复形选读}\fancyhead[R]{\small\shorttitle}\fancyfoot[C]{\thepage}
\newcommand{\sourcepage}[1]{\par\medskip\noindent{\color{gray}\small〔原文第 #1 页〕}\par\smallskip}
\newcommand{\entry}[2]{\par\medskip\noindent\textbf{#1\quad #2}\quad}
\newenvironment{proofcn}{\par\noindent\textbf{证明。}}{\hfill$\square$\par}
\newcommand{\tnote}[1]{\par\smallskip\noindent{\small\color{blue!45!black}\textbf{译注。}#1}\par\smallskip}
\newcommand{\omitsection}[2]{\section*{#1\quad #2}\addcontentsline{toc}{section}{#1\quad #2（留空）}\noindent 本选读稿不译本节；保留节号与标题作为位置标记。\par\vspace{7mm}}
\newcommand{\B}{\mathcal B}\newcommand{\C}{\mathcal C}
\newcommand{\IP}{\operatorname{IP}}\newcommand{\IA}{\operatorname{IA}}\newcommand{\EP}{\operatorname{EP}}\newcommand{\EA}{\operatorname{EA}}
\newcommand{\rk}{\operatorname{rk}}\newcommand{\lk}{\operatorname{lk}}\newcommand{\del}{\operatorname{del}}\newcommand{\con}{\operatorname{con}}
\newcommand{\lex}{\mathrm{lex}}\newcommand{\Gale}{\mathrm{Gale}}
\newcommand{\Q}{\mathbb Q}\newcommand{\R}{\mathbb R}\newcommand{\Z}{\mathbb Z}
\newcommand{\QI}{\mathrm{QI}}\newcommand{\QE}{\mathrm{QE}}\newcommand{\QC}{\mathrm{QC}}
\newcommand{\cm}{\mathrm{CM}}
\newcommand{\real}[1]{\lVert#1\rVert}
\newcommand{\blankrefs}{\section*{参考文献（留空）}\addcontentsline{toc}{section}{参考文献（留空）}正文保留原文的文献引注；完整书目信息请查本目录中的原文 PDF。本选读稿不翻译文末参考文献表。}
\hypersetup{pdfauthor={中文选读译稿}}
\begin{document}

\begin{titlepage}\centering
\vspace*{20mm}{\LARGE 有序单纯复形的拟拟阵类\par}
\vspace{6mm}{\large Quasi-matroidal classes of ordered simplicial complexes\par}
\vspace{10mm}{\Large José Alejandro Samper\par}
\vspace{4mm}华盛顿大学数学系，美国西雅图\par
\vspace{12mm}{\Large 中文分节选译\par}
\vspace{10mm}\begin{minipage}{.87\textwidth}
原文：用户提供的 arXiv:1608.04025 PDF，共 34 页；首页标注日期为 2018 年 5 月 18 日，侧栏标注 arXiv v1（2016 年 8 月 13 日）。本稿以该 PDF 为准，不替换为其他版本。

\textbf{翻译范围：}摘要；第 1、2、3 节；第 4 节开头（含定义 4.1、定理 4.2、4.3、问题 4.4）；第 4.1、4.2 小节。其余分节保留原编号、标题和留空说明。所选部分包括原文的论证和证明，不是内容提要。

\textbf{建议读法：}先读第 2 节的记号，再读第 3 节、第 4 节及 4.1、4.2。第 1 节保留了作者的研究背景和动机；其中后半部分的 $h$ 向量与猜想可在初读时跳过。

\textbf{记号：}本篇把空集算作面，面的秩是其元素个数；秩为 $d$ 的复形维数为 $d-1$。为避免源稿中删除、收缩符号混淆，译稿分别写作 $\del_e\Psi$、$\con_I\Psi$。Gale 偏序统一用 $\le_{\Gale}$ 表示。

\textbf{版本说明：}原稿有少量明显笔误，译稿在对应位置用“译注”说明；保留原定理编号及文献引注。文中对研究进展的叙述反映原稿写作时间。
\end{minipage}\vfill 2026 年 10 月\par
\end{titlepage}
\tableofcontents\newpage
\section*{摘要}\addcontentsline{toc}{section}{摘要}
\sourcepage{1}
我们引入有序单纯复形的\textbf{拟拟阵类}（quasi-matroidal class）这一概念：在有序单纯复形的范围内，它近似地体现了拟阵的等价公理刻画（cryptomorphism）。拟拟阵类包含纯移位单纯复形和有序拟阵的独立复形。它的关键性质是：如果一个固定单纯复形对顶点集的每一种排序都属于该类，那么它必为拟阵的独立复形。此类对象的一些例子已隐含地出现在拟阵论的文献中。我们引入若干这样的类，突出拟阵理论的不同方面，以及它与移位单纯复形理论的相似性。例如，我们把 Tutte 多项式和无破圈复形（nbc complex）等对象的研究推广到一个拟拟阵类，从而能为移位复形定义这些对象。此外，有些拟拟阵类适合采用不能直接用于拟阵论的归纳方法。作为例子，我们给出一个适当框架，在其中重新表述并推广 Stanley 关于拟阵 $h$ 向量的猜想；预计在这个框架中，可以使用仅研究拟阵时无法使用的方法处理它。这个新猜想对纯移位单纯复形以及秩不超过 4 的拟阵成立。

作者感谢 Isabella Novik 提供由 NSF 项目 DMS-1361423 资助的研究助理职位。

\section{引言}
\sourcepage{2}
cryptomorphism 是 Birkhoff [Bir67] 提出的一个非正式数学概念，旨在描述这样一种现象：某类对象可以用若干并非显然等价的方式来描述。拟阵是研究数学中独立性概念的一种抽象工具，可以由种类繁多、彼此等价而风格不同的公理来定义。经典的拟阵等价公理包括独立性公理、回路公理、基交换公理、次模性公理、闭集交换公理及闭包公理等。每组公理都提供了一种研究拟阵的自然途径。此外，拟阵论中许多定理似乎与某些特定公理有深刻联系：从一种角度很容易证明，从另一种角度却很困难。关于拟阵论及现有的许多等价刻画，可参阅 Oxley [Oxl92]、Welsh [Wel76] 的书，以及 Björner [Bjö92] 和 Ardila [Ard15] 的书中章节。

多年来，还出现了多种其他等价刻画，它们也对许多其他目的有用。有趣的例子包括：单纯复形理论中诱导子复形的纯性；交换代数中独立复形及其全部诱导子复形的 Stanley--Reisner 环均为 Cohen--Macaulay 环的性质（见 [Sta96]）；多胞形理论中的拟阵基多胞形（见 [GGMS87]）；以及优化中贪心算法对变化的权值均有效的性质（见 [Bjö92]）。

拟阵的许多定理似乎都附带着一条公理或一组自然的公理，也就是说，这些公理在证明目标性质时起关键作用。例如，拟阵可壳这一事实自然地来自交换公理，而内部活动性理论来自可壳性。另一方面，用诱导子复形纯性刻画拟阵，似乎是独立性公理的自然结果。无破圈复形的行为以及外部活动性理论，则似乎由回路公理支配。沿着这一启发式思路，Tutte 多项式的行为应由交换公理和回路公理共同控制，因为其内部与外部活动性有自然解释。

有两种特别有趣的等价刻画来自有序拟阵理论。Björner [Bjö92] 证明：一个单纯复形是拟阵的独立复形，当且仅当对顶点集的每一种排序，由此诱导的极大面的字典序都是壳序。另一个出色的刻画由 Gale [Gal68] 给出，它是逐坐标序中的极小性质；这种序现在称为 Gale 序。固定集合 $E$ 的一族 $d$ 元子集是某个拟阵的基族，当且仅当对 $E$ 的每一种排序，字典序最小的极大面也逐坐标最小。也就是说，若 $b_1<\cdots<b_r$ 是字典序最小基的元素，而 $b'_1<\cdots<b'_r$ 是任意其他极大面的元素，那么对每个 $i$ 都有 $b_i\le b'_i$。

\sourcepage{3}
后两种刻画特别有趣，是因为它们给出了有序单纯复形的一种性质，并要求它对底集的所有可能排序都成立。在拟阵论中，如同在线性代数中处理一族向量时一样，为独立复形的底集赋予一个全序是常见做法。例如，拟阵中得到广泛研究的无破圈复形（见 [Bry77]），只有固定底集的一个序后才能构造。不同的底集排序事实上可能产生许多互不同构的无破圈复形。另一个例子来自 Tutte 多项式理论（见 [Tut54, Cra69]）：每个拟阵都可以对应一个整系数二元多项式。Tutte 多项式编码了所有满足线性删除--收缩递推式的拟阵不变量。已知它的系数为非负整数；但只有固定底集的一个序后，才有其系数的组合解释。

更仔细研究底集排序的另一动机来自移位复形理论。移位复形是一个引人注目的单纯复形类，其结构简单、优雅且实用，因而受到重视。它们出现在 Kruskal--Katona 单纯复形面计数定理的证明中（见 [Kru63, Kat68]），也出现在 Kalai 的代数移位理论 [Kal02] 中；后者在进行同样的面计数时，还能跟踪原复形的更精细不变量。移位单纯复形的定义是：有序复形中的大顶点可以换成小顶点，而所得集合仍在复形中。这与通过交换公理描述的拟阵论颇为相似。不过，两类复形相当不同：前者依赖底集的特定排序，且包含许多非拟阵复形；而固定底集上的大多数拟阵，无论怎样排序都不是移位复形。

然而，两类对象之间的相似性十分突出。例如，在移位复形一侧假设纯性后，两类都具有相当自然的壳序（拟阵须先排序），而由这些壳序读出的组合不变量的行为也非常相似。此外，两类都有很灵活的限制与收缩理论，都对某种对偶封闭，其相应的 Gale 序也都有最小元。两类的交集同样值得注意：既为移位复形、又为拟阵独立复形的有序复形有时称为 Schubert 拟阵；它们对应于带标架的 Grassmann 流形的 Schubert 分层中一般位置点所确定的拟阵。

\sourcepage{4}
更为突出而神秘的相似性，来自 Eckmann [Eck45] 和 Friedman [Fri98] 引入的组合 Laplace 算子理论。对于单纯复形 $\Delta$，设 $(C_\bullet(\Delta),\partial)$ 是其在 $\R$ 上的单纯链复形；使用各维链群的自然面基，得到对偶复形 $(C^\bullet(\Delta),\delta)$。对每个整数 $k$，算子
\[
D_k:=\delta\partial+\partial\delta
\]
称为 $\Delta$ 的 Laplace 算子；它是 $C_k(\Delta)$ 上的自伴算子，特征值均为非负实数。因此，人们希望像谱图论那样，把 $D_k$ 的谱理论与 $\Delta$ 的组合结构联系起来：图可以视为一维单纯复形，经典的谱理论是这里的一个特例。

一系列论文 [KRS00, Den01, DR02, Koo04, Duv05] 证明，拟阵独立复形和移位单纯复形的 Laplace 算子的特征值都是整数。此外，可以将特征值放入一个二元生成函数中，称为\textbf{谱多项式}。它满足一种特殊递推式，类似于拟阵的删除--收缩递推式，但多了一个来自相对拓扑的误差修正项。这是很罕见的性质：固定顶点集上的大多数单纯复形，其 Laplace 算子并没有整谱。这自然引出以下反复被提出的问题。

\entry{问题 1.1}{[Rei01, Duv14, DR02]}是否存在一个单纯复形类，既包含拟阵独立复形，又包含移位单纯复形，并能解释整 Laplace 谱现象？

更仔细研究有序复形的另一理由来自拟阵的 $f$ 向量理论。秩为 $d$ 的单纯复形 $\Delta$ 的 $f$ 向量 $(f_0,\ldots,f_d)$ 计数各秩的面：$f_i$ 是秩为 $i$ 的独立集（即独立复形的面）的数目。自然的问题是刻画拟阵可能具有哪些 $f$ 向量。对另一些复形类，这个问题已得到完整而成功的回答，例如全体单纯复形 [Kru63, Kat68]、Cohen--Macaulay 单纯复形 [Sta77]、单纯多胞形 [BL81, Sta80]。拟阵的 $h$ 向量是 $f$ 向量的一种可逆变换，有时更方便。因此，等价的问题是分类拟阵可能具有的 $h$ 向量。这里的优点在于：固定底集的一个排序后，拟阵各基的字典序壳序给出了 $h$ 向量理论的组合实现，详见 Björner [Bjö92]。

虽然人们认为拟阵的 $h$ 向量族相当复杂，完整分类希望渺茫，但其可能取值仍受到若干限制。Adiprasito、Huh、Katz [AHK15] 在 Huh 与 Katz 的先前工作 [Huh12, HK12] 基础上，证明了一个惊人的结果：拟阵无破圈复形的 $f$ 向量是对数凹的，从而解决了 Heron、Rota 和 Welsh 的一个长期猜想。这对拟阵的 $f$ 向量族施加了很强的限制。

\sourcepage{5}
拟阵论中最引人关注的问题之一，是拟阵独立复形的 $h$ 向量。对一个有序的秩 $d$ 拟阵，基的字典序是独立复形的壳序；无破圈基也有同样性质。于是，拟阵的独立复形和无破圈复形都是 Cohen--Macaulay 的，相应的 $h$ 向量因而都是 $O$ 序列。换言之，存在一族对取因子封闭的单项式 $O$，使其对每个 $i=1,\ldots,d$ 恰有 $h_i$ 个 $i$ 次单项式。已知的一般构造并非组合构造，且所得 $O$ 的结构与拟阵本身的结构没有多少联系。容易找到一些不是拟阵 $h$ 向量的 $O$ 序列，因此可以问：拟阵的 $h$ 向量是否还须满足其他条件？Stanley 于 1977 年提出以下猜想。

\entry{猜想 1.2}{[Sta77]}拟阵独立复形的 $h$ 向量是纯 $O$ 序列。

纯 $O$ 序列的含义是：存在实现拟阵独立复形 $h$ 向量的多重复形 $O$，并且 $O$ 中关于整除关系极大的单项式全都具有相同次数。过去几十年，Stanley 猜想受到广泛关注，已知对多种特殊拟阵类成立 [Mer01, Sch10, Oh13, MNRIVF12, DLKK12, HSZ13, CKV14, Oh11, KS15, Dal15]。关于当时的进展，可参阅 [KS15] 或 Dall [Dal15]。

关于拟阵 $h$ 向量，还知道更多结果。Hibi [Hib89] 找到一组纯 $O$ 序列满足的不等式；Chari [Cha97] 给出拟阵独立复形的一种拓扑分解，由此推出其 $h$ 向量满足 Hibi 不等式。Swartz [Swa03] 又在独立复形的 Stanley--Reisner 环的 Artin 化约中，给出这些不等式的代数版本。Juhnke-Kubitzke 和 Van Dinh [JKD16] 基于 Huh [Huh15] 的工作，证明可表示拟阵的无破圈复形的 $h$ 向量也满足这类不等式。

在 [KS15] 中，Klee 与本文作者提出了 Stanley 猜想的一个更精细版本：它预测存在一个多重复形 $O$，其组合结构与原拟阵的组合结构相关。思路是利用壳序：$O$ 中的每个单项式对应于拟阵的一个基，并取决于该基的限制集。这条路线有两个主要障碍。首先，沿壳序逐步构造拟阵时，中间复形并不一定来自拟阵。其次，不能指望整个过程中都保持纯性，因此归纳框架中需要一种替代纯性的条件。

本文主要讨论三个拟拟阵类。它们是有序单纯复形的类，包含所有有序拟阵和所有纯移位复形，并且如果固定单纯复形对于顶点集的每种排序都属于指定的拟拟阵类，那么它必是拟阵的独立复形。

\sourcepage{6}
文献中已隐含地存在拟拟阵类的例子。例如，极大面的字典序为壳序的有序复形组成这样一个类；Gale 序有最小元的有序复形组成另一个例子。本文将描述若干新的拟拟阵类，其中三个类与经典的三种等价公理紧密相关：独立性公理、交换公理、回路公理。三个类两两不同，分别具有拟阵的一些有趣性质，并澄清拟阵与移位复形之间的相似性。

我们关于拟拟阵类的结果概括如下。
\begin{itemize}
\item 每个类都承载一部分拟阵理论，从而提供一种分类拟阵性质的方式：按照为了得到类似结论而需要推广哪些经典拟阵性质进行分类。
\item 独立性拟拟阵类保证许多诱导子复形是纯的，并正式说明独立性公理与诱导子复形纯性之间的依赖关系。
\item 交换拟拟阵类是按字典序可壳的复形类的子类。它有有意义的内部活动性理论，并对所谓 Gale 截断封闭，适合对极大面数作归纳。
\item 回路拟拟阵类具有良好的基本回路理论及无破圈复形理论。
\item 交换类与回路类的交集保证无破圈复形纯且可壳，并存在行为良好的 Tutte 多项式，其系数可用内部与外部活动性作组合解释。
\item 同属独立性类和交换类、且满足另一技术条件的复形，容许重新表述 Klee 与作者的猜想 [KS15]。这一猜想将 Stanley 猜想中的纯性要求转化为偏序条件，适合对基数（即极大面数）归纳。
\item 新猜想对秩不超过 4 的拟阵的 Gale 截断，以及 Dall [Dal15] 定义的内部完美拟阵的 Gale 截断成立。此外，要验证秩 $d$ 复形的猜想，只须验证顶点数不超过 $2d-1$ 的复形。
\item 还证明新猜想对纯移位复形成立，从而给出 Schubert 拟阵满足 Stanley 猜想的新证明。
\end{itemize}

\sourcepage{7}
本文安排如下：第 2 节给出背景及定义；第 3 节引入拟拟阵类，并研究一些基本性质；第 4 节讨论独立性、交换和回路拟拟阵类及各自的基本性质；第 5 节研究初始极大面性质，它是有序复形局部结构上的一个条件，使我们能灵活运用保持某些拟拟阵类的组合运算；第 6 节研究交换类与回路类的交集，并发展 Tutte 多项式和无破圈复形理论；第 7 节对具有初始极大面性质、且同属独立性类与交换类的复形，放宽 Klee 与作者的猜想，并在最后对移位复形证明新猜想；第 8 节给出公开问题、今后研究计划的简述，以及与现有文献联系的一些评论。

\textbf{致谢。}作者感谢 Isabella Novik、Steve Klee、Vic Reiner、Ernest Chong、Jeremy Martin 和 Guillaume Chapuy 提供的许多有趣讨论与建议。
\section{预备知识}
\textbf{有序单纯复形}是二元组 $\Psi=(E,\Delta)$，其中 $E$ 是有限全序集，$\Delta\subseteq2^E$ 是单纯复形，即若 $A\in\Delta$ 且 $B\subset A$，则 $B\in\Delta$。全文采用拟阵术语：$\Delta$ 的元素称为\textbf{独立集}；按包含关系极大的独立集称为\textbf{基}，基族记作 $\B$。若所有基的基数相同，称 $\Psi$ 为\textbf{纯的}。字典序最小基记为 $B_0$。不属于 $\Delta$ 的极小集合称为\textbf{回路}，回路族记作 $\C$。独立集的\textbf{秩}等于其基数；$A\subseteq E$ 的秩是包含于 $A$ 的独立集的最大秩。稍微滥用记号，将 $\Psi$ 的秩定义为 $E$ 的秩，通常记作 $d$。两个有序复形 $\Psi=(E,\Delta)$、$\Psi'=(E',\Delta')$ 称为\textbf{同构}，若 $|E|=|E'|$，且 $E$ 与 $E'$ 之间唯一的保序双射诱导 $\Delta$ 与 $\Delta'$ 之间的双射。

$\Psi$ 的\textbf{环}是 $E$ 中不属于任何基的元素。不是环的元素称为\textbf{顶点}；顶点集记为 $V(\Psi)$。属于每个基的元素称为\textbf{余环}。对 $A\subseteq E$，\textbf{限制}定义为
\[
\Psi|_A=(A,\Delta|_A),\qquad \Delta|_A=\{I\in\Delta:I\subseteq A\}.
\]
不是余环的元素 $e$ 的\textbf{删除}，是限制到 $E\setminus\{e\}$ 所得复形。不是环的元素 $e$ 的\textbf{收缩}为
\[
\con_e\Psi=(E\setminus\{e\},\operatorname{Link}_\Delta(e)),\qquad
\operatorname{Link}_\Delta(e)=\{I\in\Delta:e\notin I,\ I\cup\{e\}\in\Delta\}.
\]
独立集 $I$ 的收缩 $\con_I\Psi$ 是按任意顺序收缩 $I$ 中顶点所得的复形。令 $B_{I,0}$ 表示该收缩的字典序最小基。
\tnote{源稿第 7 页将删除记作 $\Psi/\{e\}$、收缩记作 $\Psi\setminus\{e\}$，但后面部分定理反过来使用这两个符号。本稿统一采用 $\del_e\Psi$ 与 $\con_I\Psi$，以运算的文字定义为准。这里的收缩就是取链接（link），删除就是取诱导子复形。}

\sourcepage{8}
复形 $(\{e\},\{\varnothing\})$ 记作 $\Psi_{\rm loop}$，复形 $(\{e\},\{\varnothing,\{e\}\})$ 记作 $\Psi_{\rm coloop}$。

给定 $\Psi=(E,\Delta)$、$\Psi'=(E',\Delta')$，$E$ 和 $E'$ 的一个\textbf{交错排序} $s(E,E')$ 是全序集 $X$，附有保序嵌入 $j:E\to X$、$j':E'\to X$，满足 $j(E)\cap j'(E')=\varnothing$ 且 $j(E)\cup j'(E')=X$。相应的\textbf{连接}（join）
\[
\Psi*_{s(E,E')}\Psi':=(s(E,E'),\Delta*\Delta')
\]
的独立集是形如 $j(I)\cup j'(I')$ 的集合，其中 $I\in\Delta$、$I'\in\Delta'$。若两个复形秩相同，则\textbf{连通和} $\Psi\#\Psi'$ 是通过唯一的保序双射将 $B_0$ 与 $B'_0$ 识别所得到的复形。$\Psi$ 的\textbf{秩 $k$ 骨架} $\Psi^{(k)}=(E,\Delta^{(k)})$ 由 $\Delta$ 中秩不超过 $k$ 的独立集构成。
\tnote{这里的“连通和”按作者给出的识别构造理解；不要直接套用拓扑流形中“挖球后沿边界粘合”的定义。秩 $k$ 骨架对应通常的维数 $k-1$ 骨架。}

设 $B$ 是基。若 $e\in E\setminus B$，且存在回路 $C\subseteq B\cup\{e\}$ 使 $e=\min C$，则称 $e$ 关于 $B$ \textbf{外部活动}；不外部活动时称为\textbf{外部被动}。外部活动和外部被动元素集分别记为 $\EA(B)$、$\EP(B)$。若 $b\in B$，且 $B$ 是包含 $B\setminus\{b\}$ 的字典序最小基，则称 $b$ \textbf{内部活动}。等价地，不存在 $b'<b$、$b'\notin B$ 使
\[
(B\setminus\{b\})\cup\{b'\}\in\B.
\]
非内部活动的 $b\in B$ 称为\textbf{内部被动}。相应集合记为 $\IA(B)$ 与 $\IP(B)$。

$\Psi$ 的\textbf{破圈}是形如 $C\setminus\{\min C\}$ 的子集，其中 $C$ 为回路。\textbf{无破圈复形} $\operatorname{nbc}(\Psi)=(E,\Gamma)$ 的基，是 $\Psi$ 中不包含破圈的那些基。

纯有序复形的\textbf{Gale 序}是基族上的偏序：若
\[
B=\{b_1<\cdots<b_d\},\qquad B'=\{b'_1<\cdots<b'_d\},
\]
则 $B\le_{\Gale}B'$ 当且仅当对每个 $i$ 都有 $b_i\le b'_i$。设 $J$ 为这个偏序集的序理想，则在 $J$ 处的\textbf{Gale 截断} $\Psi[J]=(E,\Delta[J])$ 是以 $J$ 的元素为基的复形。$\Psi$ 的\textbf{内部偏序集} $\operatorname{Int}(\Psi)$ 由关系
\[
B_1\le_{\rm int}B_2\quad\Longleftrightarrow\quad\IP(B_1)\subseteq\IP(B_2)
\]
定义。
\tnote{以上是原稿的用语。在一般有序复形上，内部被动集可能不能区分基，因此该关系未必是偏序；无破圈基也未必存在。后文在具有适当性质的类中使用这些对象。第 4.2 节的可壳性保证相应内部被动集确实给出偏序。}

有序复形 $\Psi$ 是\textbf{有序拟阵}，当且仅当下列三条等价公理中任意一条成立。
\begin{enumerate}[label=\roman*.]
\item \textbf{独立性公理：}若 $I_1,I_2$ 独立且 $|I_1|>|I_2|$，则存在 $i\in I_1\setminus I_2$，使 $I_2\cup\{i\}\in\Delta$。
\item \textbf{交换公理：}若 $B_1,B_2$ 为基，$b_1\in B_1\setminus B_2$，则存在 $b_2\in B_2\setminus B_1$，使 $(B_1\setminus\{b_1\})\cup\{b_2\}$ 为基。
\item \textbf{回路公理：}若 $C_1,C_2$ 为不同回路且 $c\in C_1\cap C_2$，则 $(C_1\cup C_2)\setminus\{c\}$ 包含一个回路 $C_3$。
\end{enumerate}
\tnote{原稿第三条漏写“不同”；若允许 $C_1=C_2$，则删除回路任一元素后的集合独立，结论不成立。上文补明这一必要条件。}

拟阵对偶性的一个简单推论（例如见 [Gal68]）是：若 $\Psi$ 为有序拟阵，其 Gale 序同时有最小元和最大元。Gale 还证明，这实际上是拟阵的刻画性质。

\sourcepage{9}
\entry{定理 2.1}{[Bjö92]}单纯复形 $\Delta$ 是拟阵的独立复形，当且仅当对其顶点集的每一种排序 $E$，$\Psi=(E,\Delta)$ 的 Gale 偏序都有最小元。

注意，拟阵的定义完全没有用到 $E$ 的排序。有序集 $E$ 上秩为 $d$ 的\textbf{有序均匀拟阵} $U_{E,d}=(E,X_d)$，以 $E$ 的全体 $d$ 元子集为基。

有序复形 $\Psi=(E,\Delta)$ 称为\textbf{移位的}，若满足：对基 $B$，若 $i<j$、$i\notin B$ 且 $j\in B$，则 $(B\setminus\{j\})\cup\{i\}$ 仍为基。等价地，$\Psi$ 是某个 $U_{E,d}$ 的 Gale 截断。移位复形的 Gale 序同构于整数分拆的 Young 格的一个序理想。
\tnote{源稿在上述替换式中误写为 $B\setminus\{i\}\cup\{j\}$；按其前提该集合仍是 $B$，不能表达移位性。本稿改为“删去 $j$，加入 $i$”。此处以基描述的移位复形及后面的均匀拟阵截断描述采用纯复形语境；一般非纯移位复形应在所有面上定义替换。}

秩 $d$ 复形的 $f$ 向量为 $(f_0,f_1,\ldots,f_d)$，其中 $f_i$ 为秩 $i$ 独立集数。空集是唯一的秩 0 独立集，因此 $f_0=1$。$h$ 向量 $(h_0,\ldots,h_d)$ 由下式定义：
\begin{equation}
h(\Psi,t):=\sum_{j=0}^d h_jt^j=\sum_{j=0}^d f_jt^j(1-t)^{d-j}.\tag{1}
\end{equation}
$h$ 向量携带与 $f$ 向量相同的信息，有时更方便，尤其在通过交换代数研究单纯复形时。
\tnote{源稿的 $f$ 向量重复写了 $f_0$，式 (1) 两边混用了 $x,t$；本稿分别整理为 $f_0,f_1,\ldots,f_d$ 和同一变量 $t$。}

纯复形的\textbf{壳序}是基的排序 $B_1,\ldots,B_k$，使对每个 $i<j$，存在 $\ell<j$ 和 $b\in B_j$ 满足
\[
B_i\cap B_j\subseteq B_\ell\cap B_j=B_j\setminus\{b\}.
\]
存在壳序的复形称为\textbf{可壳的}。每个壳序都有以下性质：对每个 $j=1,\ldots,k$，存在 $B_j$ 的唯一极小子集 $R(B_j)$，它不包含于任何先前的 $B_i$（$i<j$）。并且
\begin{equation}
h(\Psi,x)=\sum_{j=1}^k x^{|R(B_j)|}.\tag{2}
\end{equation}
已知有序拟阵和纯移位复形都是可壳的，基的字典序就是壳序。这又给出一个刻画拟阵的性质。

\entry{定理 2.2}{[Gal68]}单纯复形 $\Delta$ 是拟阵的独立复形，当且仅当对其顶点集的每一种排序 $E$，$\Psi=(E,\Delta)$ 的 Gale 偏序都有最小元。
\tnote{所提供版本中，定理 2.1 和 2.2 的陈述完全重复，尽管前后文意在分别讨论字典序可壳刻画和 Gale 序刻画。本译稿保留原稿的重复陈述及编号；字典序可壳的刻画已在引言明确叙述。不要将这里的两段相同文字理解成两个不同结论。}

\sourcepage{10}
纯有序复形 $\Psi=(E,\Delta)$ 称为\textbf{顶点可分解的}，当且仅当满足下列条件之一：
\begin{itemize}
\item $\Psi$ 恰有一个基；
\item 存在顶点 $e$，使 $\del_e\Psi$ 顶点可分解且与 $\Psi$ 同秩，同时 $\con_e\Psi$ 顶点可分解。
\end{itemize}
若 $\Psi$ 顶点可分解且不止一个基，则满足第二条的 $e$ 称为\textbf{脱落顶点}（shedding vertex）。Billera 与 Provan [PB80] 的定理表明，顶点可分解复形是可壳的。他们还证明拟阵顶点可分解，其中每个非余环顶点都是脱落顶点。纯移位复形也顶点可分解：最大顶点总可作为脱落顶点。
\tnote{单基情形直接使用第一条。原稿关于拟阵的“任意顶点”须排除余环，因删除余环会降低秩；关于移位复形的最大顶点也应在非单基情形排除余环，采用最大非余环顶点。}

对集合 $A,B$，记其\textbf{对称差}为 $A\triangle B=(A\setminus B)\cup(B\setminus A)$。考虑较小集合的子集时，为简化记号省略花括号与逗号。例如，$\{1,2,3,4,5\}$ 的子集 $\{2,4\}$ 记作 $24$。

\section{有序复形的拟拟阵类}
本节用 $\mathcal A$ 表示一个有序复形类。若对任意 $\Psi,\Psi'\in\mathcal A$ 以及底集的任意交错排序 $s$，连接 $\Psi*_s\Psi'$ 都属于 $\mathcal A$，则称 $\mathcal A$ \textbf{对连接封闭}。若每个复形删除底集的最大元素后仍在 $\mathcal A$ 中，则称其\textbf{对删除封闭}。若每个复形收缩其任意独立集后仍在 $\mathcal A$ 中，则称其\textbf{对收缩封闭}。下面的概念概括了本文研究的核心对象。

\entry{定义 3.1}{}有序单纯复形类 $\mathcal A$ 称为\textbf{拟拟阵类}，若满足：
\begin{enumerate}
\item 每个有序拟阵都是 $\mathcal A$ 的对象。
\item 若单纯复形 $\Delta$ 的顶点集为 $X$，且对 $X$ 的每个全序 $E$，$(E,\Delta)$ 均属于 $\mathcal A$，则 $\Delta$ 是拟阵的独立复形。
\item 每个纯移位复形均在 $\mathcal A$ 中。
\item $\mathcal A$ 对连接、删除、收缩封闭。
\end{enumerate}

定理 2.1 和 2.2 给出了两种不同的拟拟阵类。
\entry{例 3.2}{}以下两个类是拟拟阵类：
\begin{enumerate}
\item $\mathrm{LEX}$：对连接、删除和收缩封闭，且基的字典序为壳序的纯有序复形类。
\item $\mathrm{GALE}$：对连接、删除和收缩封闭，且 Gale 偏序有唯一极小基的纯有序复形类。
\end{enumerate}

\sourcepage{11}
稍大一些、包含上述两类的例子如下。
\entry{例 3.3}{}用 $\mathrm{PURE}$ 表示满足以下递归描述的纯有序复形类：$\Psi=(E,\Delta)$ 属于 $\mathrm{PURE}$，当且仅当以下之一成立：
\begin{itemize}
\item $\Psi$ 恰有一个基；或
\item 若 $v$ 是 $\Psi$ 的最大非余环顶点，则 $\del_v\Psi$ 属于 $\mathrm{PURE}$，且与 $\Psi$ 同秩；此外，$\Psi$ 的每个收缩都属于 $\mathrm{PURE}$。
\end{itemize}
容易看出纯移位复形和拟阵都属于 $\mathrm{PURE}$。事实上，由以下经典定理，$\mathrm{PURE}$ 是拟拟阵类。

\entry{定理 3.4}{[Sta96]}单纯复形 $\Delta$ 是拟阵的独立复形，当且仅当其每个诱导子复形都是纯的。

$\mathrm{PURE}$ 解释了纯移位复形和拟阵之间最初几种良好的相似性。
\entry{定理 3.5}{}$\mathrm{PURE}$ 中的每个有序复形都顶点可分解，因而可壳。

虽然以上例子解释了移位复形和拟阵独立复形的一些相似性，但这些类太大，包含很多远离二者的复形。例如，它们都包含以 $12,13,24$ 为基的复形 $\Psi=([4],\Delta)$。这个复形是三条边组成的路径，也是 Laplace 算子具有非整数谱的典型例子。因此，希望研究更小的拟拟阵类，使其成员与拟阵和移位复形共享更深的结构性质。

为此，引入拟拟阵类的几种基本构造。给定拟拟阵类 $\mathcal A,\mathcal A'$，由同时属于二者的复形组成的交集 $\mathcal A\cap\mathcal A'$ 显然也是拟拟阵类。若 $\mathcal A'$ 的全部元素都属于 $\mathcal A$，则称 $\mathcal A'$ 为 $\mathcal A$ 的\textbf{子类}。

\section{三个拟拟阵类}
本节引入三个新的拟拟阵类，全文将围绕它们展开研究。
\sourcepage{12}
\entry{定义 4.1}{}设 $\Psi=(E,\Delta)$ 是有序单纯复形。用以下公理定义类 $\QI,\QE,\QC$。
\begin{itemize}
\item \textbf{拟独立性公理（QI）：}$\Psi$ 是纯的；对任意独立集 $I_1,I_2$，若 $|I_1|>|I_2|$，并存在 $I\subseteq I_1\cap I_2$ 使 $I_1\setminus I\subseteq B_{I,0}$，则存在 $e\in I_1\setminus I_2$，使 $I_2\cup\{e\}$ 独立。
\item \textbf{拟交换公理（QE）：}$\Psi$ 是纯的；对任意基 $B_1,B_2$，若 $b_1\in B_1\setminus B_2$ 且 $b_1>\max(B_2\setminus B_1)$，则存在 $b_2\in B_2\setminus B_1$，使 $(B_1\setminus\{b_1\})\cup\{b_2\}\in\Delta$。
\item \textbf{拟回路公理（QC）：}若 $C_1,C_2$ 为不同回路，$c\in C_1\cap C_2$ 且 $c<\max(C_1\triangle C_2)$，则 $(C_1\cup C_2)\setminus\{c\}$ 包含一个回路 $C_3$。
\end{itemize}

首先要说明 $\QI,\QE,\QC$ 是拟拟阵类。容易看出，有序拟阵属于这三个类。拟拟阵类的第二条公理也容易验证：去掉 QI、QE、QC 中关于排序的限制，就得到经典的独立性、交换和回路公理。另一方面，证明移位单纯复形满足这些公理是一道有趣的练习。

\entry{定理 4.2}{}若 $\Psi$ 是纯移位复形，则 $\Psi$ 属于 $\QI,\QE,\QC$。
\begin{proofcn}
为证 QI，注意若独立集 $I$ 与 $B_0$ 不交，则 $B_{I,0}$ 是 $B_0$ 的前 $d-|I|$ 个元素构成的初始子集，其中 $d=|B_0|$。设独立集 $I_1,I_2$ 满足 QI 的条件，其见证集为 $I\subseteq I_1\cap I_2$。于是 $I_1\setminus I$ 包含于 $B_{I,0}$，而 $B_{I_2,0}\subseteq B_{I,0}$。必有
\[
B_{I_2,0}\cap(I_1\setminus I)\ne\varnothing.
\]
否则，因为二者都是 $B_{I,0}$ 的子集，将有
\[
\big|B_{I_2,0}\cup(I_1\setminus I)\big|
=(d-|I_2|)+(|I_1|-|I|)>d-|I|=|B_{I,0}|,
\]
矛盾。交集中的元素属于 $I_1\setminus I_2$，并可加入 $I_2$ 保持独立。

为证 QE，若基 $B_1,B_2$ 中的 $b_1\in B_1\setminus B_2$ 大于 $B_2\setminus B_1$ 的全部元素，移位性允许用其中任意 $b_2$ 替换 $b_1$。

为证 QC，取满足公理条件的回路 $C_1,C_2$ 和公共元素 $c$。假设 $I=(C_1\cup C_2)\setminus\{c\}$ 独立，并令 $c_1=\max(C_1\triangle C_2)$。移位性保证
\[
I'=(I\setminus\{c_1\})\cup\{c\}
\]
独立，但它包含 $C_1$ 或 $C_2$ 中的一个，矛盾。
\end{proofcn}
\tnote{原定理省略“纯”，但 QI、QE 的定义要求纯性，故在本节语境中须加上它。QI 证明中原稿先在 $I\cap B_0=\varnothing$ 的特例下描述 $B_{I,0}$，随后直接使用；一般情形应将“$B_0$ 的前 $d-|I|$ 个元素”改读为“$B_0\setminus I$ 的前 $d-|I|$ 个元素”。由此仍有 $B_{I_2,0}\subseteq B_{I,0}$，后面的计数证明不变。}

显然，每个拟拟阵类都能转化为一种拟阵的等价公理刻画。虽然拟阵的各种等价公理产生相同的对象类，但拟拟阵类却有多种，各自突出拟阵理论的不同方面。下一个定理说明，刚定义的三个类确实不同。

\entry{定理 4.3}{}$\QI,\QE,\QC$ 两两不同，并且任何一类都不包含在另一类中。
\sourcepage{13}
\begin{proofcn}
对每一对公理，只须给出满足其中一个而不满足另一个的例子。
\begin{itemize}
\item $\Psi_1=(\{1,2,3,4\},\Delta_1)$ 的基为 $12,13,14,34$。它满足 QI、QE，但回路 $23,24$ 违反 QC。
\item $\Psi_2=(\{1,2,3,4\},\Delta_2)$ 的基为 $14,24,23,34$。它满足 QI；基 $14,23$ 表明它不满足 QE，回路 $12,13$ 表明它不满足 QC。
\item $\Psi_3=(\{1,2,3,4\},\Delta_3)$ 的基为 $12,13,23,34$。它满足 QC，但不满足 QI、QE。
\item $\Psi_4=(\{1,2,3,4,5\},\Delta_4)$ 的基为 $13,14,23,24,25$。它满足 QE，但不满足 QI、QC。
\end{itemize}
\end{proofcn}
\tnote{第二个例子中，原稿将用于反驳 QC 的回路写成 $13,14$；但 $14$ 本来就是基，不可能是回路。应使用 $12,13$：公共元素 $1<\max\{2,3\}$，而删除公共元素后的 $23$ 是基，不含回路。}

这个证明表明 $\QC$ 不包含于 $\QI\cap\QE$。另一方面，一个直接的图论练习表明，属于 $\QE\cap\QC$ 的秩 2 复形也属于 $\QI$。不过，这个证明在秩 3 时失效，因此自然提出如下问题。

\entry{问题 4.4}{}$\QI\cap\QE$、$\QI\cap\QC$、$\QE\cap\QC$ 与 $\QI\cap\QE\cap\QC$ 是否全都不同？

不论答案如何，这些定义给出了若干单纯复形类；它们各自与有序拟阵类共享某些结构性质，并不奇怪。本节余下部分的目标是为这些类建立理论的第一步：既给出拟阵性质的某些类似结论，也给出可以在指定类内部进行的多种构造。

\subsection{拟独立性类}
现在说明 $\QI$ 是具有额外结构的拟拟阵类。
\entry{定理 4.5}{}设 $\Psi=(E,\Delta)$、$\Psi'=(E',\Delta')$ 属于 $\QI$。
\begin{enumerate}[label=\roman*.]
\item 对 $E,E'$ 的任意交错排序 $s(E,E')$，连接 $\Psi*_{s(E,E')}\Psi'$ 属于 $\QI$。
\item 若 $v\in E\setminus B_0$ 是顶点，则收缩 $\con_v\Psi$ 属于 $\QI$。
\item 若 $v$ 是 $\Psi$ 的最大非余环顶点，则 $\del_v\Psi$ 属于 $\QI$。
\item 若 $0\le k\le\rk(\Psi)$，则秩 $k$ 骨架 $\operatorname{Skel}_k(\Psi)$ 属于 $\QI$。
\item 若 $\rk(\Psi)=\rk(\Psi')$，则连通和 $\Psi\#_\phi\Psi'$ 属于 $\QI$。
\end{enumerate}
第 i、ii、iii 项表明，$\QI$ 是 $\mathrm{PURE}$ 的拟拟阵子类。
\tnote{本稿保留原稿第 ii 项的限制 $v\notin B_0$ 及其后这句推论。按定义 3.1，“对收缩封闭”要求任意独立集的收缩；因此仅凭列出的这三项还不能直接得到该推论，原稿此处的论述有省略。环不能收缩，故第 ii 项明确限于顶点。}

\sourcepage{14}
\begin{proofcn}
第 i 项来自以下事实：连接保持纯性，并与链接运算交换。即若 $I_1,I_2$ 分别是 $\Delta_1,\Delta_2$ 的面，则
\[
\operatorname{Link}_{\Delta_1*\Delta_2}(I_1\cup I_2)
=\operatorname{Link}_{\Delta_1}(I_1)*\operatorname{Link}_{\Delta_2}(I_2).
\]
对于第 ii 项，若 $I_1,I_2$ 以及 $I\subseteq I_1\cap I_2$ 在 $\operatorname{Link}_\Delta(v)$ 中满足公理假设，则 $I_1\cup\{v\}$、$I_2\cup\{v\}$、$I\cup\{v\}$ 在 $\Psi$ 中满足相应假设。因此可在 $\Psi$ 中使用 QI 公理得到结论。

为证第 iii 项，先说明：若 $I$ 是不含 $v$ 的独立集，则 $v\notin B_{I,0}$。否则，令
\[
B=(I\cup B_{I,0})\setminus\{v\}.
\]
对 $B_0$ 与 $B$ 使用 QI，得到 $u\in B_0\setminus B$ 使 $B\cup\{u\}$ 独立。$u$ 不在基 $I\cup B_{I,0}$ 中，故不是余环，因而 $u<v$。于是 $(B\cup\{u\})\setminus I$ 是收缩 $\con_I\Psi$ 的一个基，且字典序小于 $B_{I,0}$，矛盾。这也使删除后的独立集及其字典序最小补全可沿用原复形的 QI 公理。

第 iv 项来自恒等式
\[
\operatorname{Link}_{\operatorname{Skel}_k\Delta}(I)
=\operatorname{Skel}_{k-|I|}(\operatorname{Link}_\Delta I),
\]
以及 $\operatorname{Skel}_k\Psi$ 的字典序最小极大面正是 $B_0$ 的字典序最小 $k$ 元面这一事实。

对于第 v 项，假设 $I_1,I_2$ 及 $I\subseteq I_1\cap I_2$ 在连通和中满足公理假设。若 $I_1,I_2$ 分别来自 $\Psi,\Psi'$，则在自然识别下 $I_1\cap I_2\subseteq B_0=B'_0$，因此
\[
B_{I,0}=B_0\setminus I=B'_0\setminus I.
\]
于是 $I_1\subseteq B_0=B'_0$，可在 $\Psi'$ 中应用公理。若两个独立集都来自同一侧，便在该侧应用 QI。
\end{proofcn}

下一个定理从特定意义上说明 QI 细化了 PURE：QI 复形中有更大一族诱导子复形为纯复形。
\entry{定理 4.6}{}设 $\Psi=(E,\Delta)\in\QI$，且 $A\subseteq E$ 满足 $\rk(A)=|B_0\cap A|$。则 $\Psi|_A$ 是纯的。
\begin{proofcn}
若 $I$ 是 $\Psi|_A$ 的独立集，可对 $B_0\cap A$ 与 $I$ 应用 QI，将 $I$ 扩充为 $\Psi|_A$ 的独立集 $B$，直到
\[
|B|=|B_0\cap A|=\rk(A).
\]
于是每个独立集都能扩充到满秩的基，故限制是纯的。
\end{proofcn}

\subsection{交换公理}
单纯复形的可壳性理论最好用极大面的语言研究，因此它自然应当由基族上的条件推出。事实表明，QE 是运用这套工具的适当对象类。

\entry{定理 4.7}{}设 $\Psi=(E,\Delta)$、$\Psi'=(E',\Delta')$ 属于 $\QE$。
\begin{enumerate}[label=\roman*.]
\item 对 $E,E'$ 的任意交错排序 $s(E,E')$，连接 $\Psi*_{s(E,E')}\Psi'$ 属于 $\QE$。
\item 若 $U\subseteq E$，且 $\Psi|_U$ 纯并与 $\Psi$ 同秩，则 $\Psi|_U\in\QE$。
\item 若 $A\subseteq E$ 包含 $B_0$，并包含所有小于 $B_0$ 中最大非余环元素的元素，则 $\Psi|_A\in\QE$。
\item 若 $v$ 是顶点，则 $\con_v\Psi\in\QE$。
\end{enumerate}
特别地，$\QE$ 是 $\mathrm{PURE}$ 的拟拟阵子类。

\sourcepage{15}
\begin{proofcn}
第 i 项来自连接保持纯性，且连接中的基对应于两侧各取一个基。因而对可以交换的每个元素，在其所属一侧应用 QE 即可。

对于第 ii 项，若 $B_1,B_2$ 是 $\Psi|_U$ 的基，对它们使用 QE，得到的仍是包含于 $U$ 的基。

对于第 iii 项，由第 ii 项，只须证 $\Psi|_A$ 纯。取其一个面 $I$，并取 $\Psi$ 中包含 $I$ 的基 $B$。对 $B,B_0$ 应用 QE，可以把 $B\setminus A$ 中的顶点逐一换成 $B_0$ 的元素，因为被换出的元素均大于可换入的非余环元素。所得基包含于 $A$，且仍包含 $I$，于是限制为纯。
\tnote{原证明误将 $I$ 写成未定义的 $F$，并写 $\min(E\setminus A)>\max B_0$。后一个不等式在 $B_0$ 含较大余环时未必成立；实际只须比较 $B_0\setminus B$ 的元素，它们都不是余环，故上述交换论证仍成立。}

对于第 iv 项，若 $B_1,B_2$ 是收缩的基，则 $B_1\cup\{v\},B_2\cup\{v\}$ 是原复形的基，可对它们应用 QE。由于 $v$ 同属于两个基，它不影响交换条件，因此收缩也满足 QE。
\end{proofcn}

从组合拓扑，特别是可壳性的角度看，QE 公理相当丰富。虽然定理 3.5 已说明 $\Psi$ 可壳，但顶点分解给出的壳序是递归得到的，其限制集因而难以研究。限制集的生成函数等于 $h$ 多项式，因此希望有一种壳序，其限制集能更直接地计算。QE 复形保证了这一点。

\entry{定理 4.8}{}若 $\Psi=(E,\Delta)\in\QE$，则 $\B$ 的字典序为壳序，且在这个壳序下 $R(B)=\IP(B)$。
\begin{proofcn}
设 $B_1<_{\lex}B_2$。目标是找到 $B_3<_{\lex}B_2$，使
\[
B_1\cap B_2\subseteq B_3\cap B_2=B_2\setminus\{e\}.
\]
对 $r-|B_1\cap B_2|$ 归纳，其中 $r$ 为复形的秩。起始情形显然成立。分两种情形。

\textbf{情形 1：}$b:=\max(B_1\triangle B_2)$ 属于 $B_2$。由 QE，存在 $b'\in B_1\setminus B_2$ 使
\[
B_3=(B_2\setminus\{b\})\cup\{b'\}
\]
为基，这个 $B_3$ 直接满足要求。

\textbf{情形 2：}$b:=\max(B_1\triangle B_2)$ 属于 $B_1$。由 QE，存在 $b'\in B_2\setminus B_1$ 使
\[
B'_1=(B_1\setminus\{b\})\cup\{b'\}
\]
为基。此时 $B_1\cap B_2\subsetneq B'_1\cap B_2$，而且 $B'_1<_{\lex}B_2$。由归纳假设，存在 $B_3$ 使
\[
B'_1\cap B_2\subseteq B_2\cap B_3=B_2\setminus\{e\},
\]
因而也满足原要求。

第二个结论直接来自限制集的定义，与 [Bjö92] 中的证明相同。
\end{proofcn}
\tnote{原稿在证明目标中将第一个交集写成 $B_1\cap B_3$；本稿依壳序定义及后续归纳步骤改为 $B_1\cap B_2$。}

\sourcepage{16}
\entry{注 4.9}{}定理 4.8 表明 $\QE\subseteq\mathrm{LEX}$，而且一般小得多。图是秩 2 单纯复形；可壳在这里就是连通。给定顶点集为 $[n]$ 的连通图 $\Delta$，若对每个 $1\le r\le n$，$\Delta|_{[r]}$ 连通，则复形 $([n],\Delta)$ 属于 LEX。然而这类图中有相当一部分不属于 QE。因此即使只比较秩 2 复形，QE 也明显小于 LEX。

下面的推论直接来自定理 4.8，推广了 Björner [Bjö92] 的结果。正如第 7 节将说明的，在这个框架中研究它会更有力。
\entry{推论 4.10}{}若 $\Psi=(E,\Delta)$ 属于 $\mathrm{LEX}$（或属于 $\QE$），则
\begin{equation}
h(\Psi,x)=\sum_{B\in\B}x^{|\IP(B)|}.\tag{3}
\end{equation}

还有一种从旧 QE 复形得到新 QE 复形的值得注意的方法，它在组合拓扑中并不常见。它表明：按定理 4.8 的壳序逐步构造 QE 复形时，部分步骤所得的复形仍属于 QE。

\entry{定理 4.11}{}若 $\Psi=(E,\Delta)\in\QE$，$J$ 是 $\operatorname{Gale}(\Psi)$ 的序理想，则 $\Psi[J]\in\QE$。此外，对 $\Psi[J]$ 的每个基 $B$，都有
\[
\IP(B,\Psi[J])=\IP(B,\Psi).
\]
\begin{proofcn}
第一个结论来自以下事实：字典序是 Gale 序的线性扩张；QE 公理给出的交换产生一个比原基更小的基，因此仍在向下封闭的 $J$ 中。

被动集的相等甚至对一般有序复形也成立，因为见证 $B$ 的内部被动性的那些基都在 Gale 序中小于 $B$：它们与 $B$ 只相差一次用小元素换大元素的替换。
\end{proofcn}
\tnote{为了保证仍属于序理想，交换后的基须在 Gale 序中更小；这里一次“换入较小元素”的交换确有这一性质。原证明第二段写成 external activity（外部活动性），但结论与所用的基替换都涉及内部被动性，本稿据此译出。}

QE 复形的一个有用性质涉及 Gale 序的结构。
\entry{引理 4.12}{}若 $\Psi=(E,\Delta)$ 满足 QE，则 $B_0$ 是 $\operatorname{Gale}(\Psi)$ 的唯一极小基。特别地，$\QE\subseteq\mathrm{GALE}$。
\begin{proofcn}
由于字典序是 Gale 序的线性扩张，$B_0$ 显然是极小元。另一方面，对任意其他基 $B$，与 $B_0$ 应用 QE 得到基 $B'$。这个交换从 $B$ 中移去一个元素，并换入更小的顶点，因此 $B'<_{\Gale}B$。
\end{proofcn}

\sourcepage{17}
另一个有用工具是 Las Vergnas 的内部偏序集。将会看到，$\operatorname{Int}(\Psi)$ 比 $\operatorname{Gale}(\Psi)$ 更粗；这一点有助于研究 Stanley 猜想。下面给出若干结构结果。

\entry{引理 4.13}{}设 $\Psi\in\QE$，$B$ 为基，且 $\IP(B)\ne\varnothing$。则存在 $b\in\IP(B)$ 和基 $B'$，使
\[
\IP(B')=\IP(B)\setminus\{b\}.
\]
特别地，$\operatorname{Int}(\Psi)$ 是分次偏序集，基的次数等于其内部被动集的基数。
\begin{proofcn}
字典序是壳序，且这个壳序的限制集恰是内部被动集，因此第一个结论来自 [Bjö92, Lemma 7.2.6]。由此立即可知，被动集的基数给出 $\operatorname{Int}(\Psi)$ 的一个分次。
\end{proofcn}
\tnote{原稿漏写 $\IP(B)\ne\varnothing$；对最小基 $B_0$，其内部被动集为空，不可能从中选出 $b$，故补明这一必要条件。}

最后验证 QE 复形的 Gale 序和内部偏序之间的以下关系。
\entry{定理 4.14}{}设 $\Psi\in\QE$，且基 $B,B'$ 满足 $\IP(B)\subseteq B'$。则 $B\le_{\Gale}B'$。因此 $\operatorname{Gale}(\Psi)$ 是 $\operatorname{Int}(\Psi)$ 的偏序扩张。
\begin{proofcn}
固定 $B$，对 $|B'\setminus B|$ 归纳。起始情形 $B=B'$ 无须证明。假设结论对所有满足 $\IP(B)\subseteq B_2$ 且 $|B_2\setminus B|<|B'\setminus B|$ 的基成立。

对 $B,B'$ 使用 QE，令 $u=\max(B\triangle B')$。由于 $\IP(B)\subseteq B\cap B'$，有 $u\notin\IP(B)$。若 $u\in B$，QE 保证存在 $b'\in B'\setminus B$，使
\[
B_2=(B\setminus\{u\})\cup\{b'\}
\]
为基。但此时 $B_2<_{\lex}B$ 且 $\IP(B)\subseteq B_2$，这不可能。因此 $u\in B'$，进而存在 $b\in B\setminus B'$，使
\[
B_2=(B'\setminus\{u\})\cup\{b\}
\]
为基。于是 $B_2<_{\Gale}B'$、$\IP(B)\subseteq B_2$，且 $|B_2\setminus B|<|B'\setminus B|$。由归纳假设，
\[
B\le_{\Gale}B_2<_{\Gale}B',
\]
即得结论。

若 $B\le_{\rm int}B'$，则 $\IP(B)\subseteq\IP(B')\subseteq B'$，故 $B\le_{\Gale}B'$。这说明 Gale 序是内部偏序的扩张。
\end{proofcn}

\omitsection{4.3}{回路公理（The Circuit axiom）}
\omitsection{5}{初始基性质（The First Basis Property）}
\omitsection{6}{Tutte 多项式和无破圈复形（Tutte polynomials and nbc complexes）}
\omitsection{7}{Stanley 猜想的一个精细化（A refinement of Stanley's conjecture）}
\omitsection{8}{问题、评论与未来方向（Questions, remarks and future directions）}
\subsection*{8.1\quad 问题（Questions；留空）}\vspace{1cm}
\subsection*{8.2\quad 评论（Remarks；留空）}\vspace{1cm}
\blankrefs
\end{document}
